% !TEX root = ../elteikthesis_en.tex
\chapter{Evaluation}
\label{ch:evaluation}

This chapter presents an empirical evaluation of the ELTE IK Assistant across four configurations, measuring answer quality, faithfulness to retrieved source material, out-of-scope refusal behaviour, retrieval effectiveness, and response latency. The goal is to quantify how much value the Retrieval-Augmented Generation component adds over a plain language model baseline, and to compare the two evaluated local models under both conditions.

% ─────────────────────────────────────────────────────────────
\section{Evaluation Setup}
\label{sec:eval-setup}

\subsection{Configurations}

Four configurations were evaluated in a factorial design crossing two models with two retrieval modes:

\begin{table}[H]
    \centering
    \begin{tabular}{|l|l|l|l|}
        \hline
        \textbf{ID} & \textbf{Model} & \textbf{RAG} & \textbf{Description} \\
        \hline\hline
        A1 & \texttt{llama3.2:3b} & No  & Baseline: model knowledge only \\
        B1 & \texttt{llama3.2:3b} & Yes & RAG with top-3 chunk retrieval \\
        A2 & \texttt{gemma3:4b}   & No  & Baseline: model knowledge only \\
        B2 & \texttt{gemma3:4b}   & Yes & RAG with top-3 chunk retrieval \\
        \hline
    \end{tabular}
    \caption{The four evaluation configurations.}
    \label{tab:eval-configs}
\end{table}

Both models were run locally via Ollama with 4-bit quantisation and temperature 0.1. No internet access was used during evaluation.

\subsection{Test Set}

The test set consists of 20 questions divided into two categories:

\begin{itemize}
    \item \textbf{15 in-scope questions} about topics present in the crawled ELTE documents: course prerequisite rules (strong and weak), Neptun registration behaviour, faculty staff information, the Computer Science BSc programme, and international partnerships. These questions have a definite correct answer that can be verified against the source documents.
    \item \textbf{5 out-of-scope questions} with no connection to the faculty corpus: general knowledge (\emph{Who won the FIFA World Cup in 2022?}), weather, cooking, geography, and humour. These questions measure whether the system correctly declines to answer rather than hallucinating.
\end{itemize}

Reference answers for in-scope questions were written manually by consulting the source documents. Each question was run through all four configurations independently; the model received no memory of previous questions.

\subsection{Metrics}

\begin{description}
    \item[ROUGE-L] The longest-common-subsequence F1 score between the generated answer and the reference answer~\cite{lin2004rouge}. Captures lexical overlap with the ground truth on a scale from 0 to 1.
    \item[Faithfulness] The proportion of words in the generated answer that appear in the retrieved chunks. A high faithfulness score indicates the model is drawing on the provided context rather than confabulating.
    \item[Refusal rate] For out-of-scope questions only: the fraction of queries for which the model produced a non-answer (e.g.\ \emph{I don't know} or \emph{This is outside my knowledge}). Higher is better for this category.
    \item[Retrieval hit-rate] For RAG configurations: the fraction of in-scope questions for which at least one of the retrieved chunks came from the expected source document. Measures the quality of the embedding-based retrieval independently of the language model.
    \item[Response time] Wall-clock latency from request to complete response, measured in milliseconds on a laptop with 8~GB RAM and no GPU.
\end{description}

% ─────────────────────────────────────────────────────────────
\section{Results}
\label{sec:eval-results}

\subsection{Overall Metric Summary}

Table~\ref{tab:eval-summary} reports the per-configuration averages across all 20 test questions.

\begin{table}[H]
    \centering
    \begin{tabular}{|l|l|c|c|c|c|r|}
        \hline
        \textbf{ID} & \textbf{Model} & \textbf{ROUGE-L} & \textbf{Faithfulness} & \textbf{Refusal} & \textbf{Hit-rate} & \textbf{Avg time} \\
        \hline\hline
        A1 & llama3.2:3b, no RAG & 0.117 & 0.377 & 0.00 & 0.667 & 47.6 s \\
        B1 & llama3.2:3b, RAG    & 0.436 & 0.711 & 0.20 & 0.667 & 70.5 s \\
        A2 & gemma3:4b,   no RAG & 0.078 & 0.355 & 0.00 & 0.667 & 121.1 s \\
        B2 & gemma3:4b,   RAG    & 0.534 & 0.844 & 0.00 & 0.667 & 61.4 s \\
        \hline
    \end{tabular}
    \caption{Evaluation results averaged over 20 test questions. Hit-rate is computed over in-scope questions only (15 questions); refusal rate is computed over out-of-scope questions only (5 questions).}
    \label{tab:eval-summary}
\end{table}

\subsection{Effect of RAG on Answer Quality}

Adding RAG context produces a consistent and substantial improvement in both lexical quality and faithfulness for both models.

\begin{itemize}
    \item \textbf{ROUGE-L} increases by a factor of 3.7$\times$ for \texttt{llama3.2:3b} (0.117 $\to$ 0.436) and by 6.9$\times$ for \texttt{gemma3:4b} (0.078 $\to$ 0.534). The larger relative gain for Gemma suggests that without retrieval it relies more heavily on parametric knowledge, which is less aligned with ELTE-specific terminology.
    \item \textbf{Faithfulness} roughly doubles in both cases (0.377 $\to$ 0.711 for Llama; 0.355 $\to$ 0.844 for Gemma), confirming that the retrieved chunks are genuinely shaping the generated text rather than being ignored.
\end{itemize}

The qualitative difference is visible in individual answers. For the question \emph{``What is a strong prerequisite?''}, the no-RAG baseline (A1) produces a generic university definition with invented ELTE-specific examples that do not appear in any source document. The RAG-augmented configuration (B1) produces an answer drawn directly from the prerequisites document: \emph{``A strong prerequisite is a prerequisite without which the follow-up subject cannot be taken in the next semester; the follow-up subject can only be registered for after the prerequisite subject has been fulfilled.''} The ROUGE-L score for this question rises from 0.13 (A1) to 0.29 (B1) and faithfulness from 0.44 to 0.80.

\subsection{Model Comparison}

Among the RAG configurations, \texttt{gemma3:4b} (B2) outperforms \texttt{llama3.2:3b} (B1) on both answer quality metrics: ROUGE-L 0.534 vs.\ 0.436, faithfulness 0.844 vs.\ 0.711. Gemma produces more concise answers that stay closer to the retrieved context, whereas Llama tends to expand its answers with surrounding commentary.

Among the no-RAG baselines, both models score similarly poorly, with Llama slightly ahead on ROUGE-L (0.117 vs.\ 0.078). Neither baseline reliably answers ELTE-specific questions correctly without retrieval.

\subsection{Retrieval Quality}

The retrieval hit-rate is 0.667 (10 out of 15 in-scope questions) and is identical across all four configurations, since the retrieval step is determined solely by the embedding model and vector index. The five missed questions share a common cause: they ask about topics (Erasmus partnerships, CEEPUS network, Dean contact details) that are present in the HTML pages but were truncated or de-prioritised during chunking because the relevant content appeared in shallow navigation-style text rather than in a dedicated content section. This is the primary bottleneck limiting further ROUGE-L and faithfulness gains: when the correct chunk is not retrieved, the language model must fall back on parametric knowledge, which tends to be inaccurate for institution-specific facts.

\subsection{Out-of-Scope Refusal}

Neither baseline model refuses any out-of-scope question: both A1 and A2 answer all five irrelevant queries confidently, generating plausible but entirely off-topic responses. Among the RAG configurations, \texttt{llama3.2:3b} with RAG (B1) refuses 1 out of 5 out-of-scope questions (20\%), triggered by the prompt instruction to answer only from provided context when no relevant chunks are retrieved. \texttt{gemma3:4b} with RAG (B2) does not refuse any out-of-scope question despite receiving the same prompt instruction, choosing instead to answer from its parametric knowledge when context is absent.

This is a notable safety gap: the RAG prompt template reduces but does not eliminate hallucination on out-of-scope inputs, and its effectiveness depends on the model.

\subsection{Response Latency}

\begin{table}[H]
    \centering
    \begin{tabular}{|l|r|r|r|}
        \hline
        \textbf{Configuration} & \textbf{Avg (s)} & \textbf{No-RAG overhead} & \textbf{Req.\ threshold} \\
        \hline\hline
        A1 — llama3.2:3b, no RAG & 47.6  & —       & exceeded \\
        B1 — llama3.2:3b, RAG    & 70.5  & +22.9 s & exceeded \\
        A2 — gemma3:4b,   no RAG & 121.1 & —       & exceeded \\
        B2 — gemma3:4b,   RAG    & 61.4  & $-$59.7 s & exceeded \\
        \hline
    \end{tabular}
    \caption{Average response latency. The non-functional requirement NF-4 sets a target of 30 seconds.}
    \label{tab:eval-latency}
\end{table}

All four configurations exceed the 30-second latency target defined in the requirements (NF-4). This is expected on CPU-only hardware with 4-bit quantised models: the bottleneck is autoregressive token generation, which scales linearly with the number of output tokens. Two observations are worth noting. First, RAG adds latency for \texttt{llama3.2:3b} (+23 s) but \emph{reduces} it for \texttt{gemma3:4b} ($-$60 s): when grounded by retrieved context, Gemma produces shorter, more focused answers, whereas without context it generates lengthy elaborations. Second, \texttt{llama3.2:3b} is substantially faster than \texttt{gemma3:4b} under both conditions, making it the more practical choice where latency matters more than marginal quality gains.

% ─────────────────────────────────────────────────────────────
\section{Discussion}
\label{sec:eval-discussion}

\subsection{RAG Effectiveness}

The evaluation results confirm that RAG is effective for this use case. Grounding the language model in retrieved faculty documents more than triples ROUGE-L scores and nearly doubles faithfulness, for both of the evaluated models. The improvements are consistent and large enough to be meaningful in practice, not merely statistical noise over a 20-question benchmark.

The fundamental mechanism is straightforward: without context, neither \texttt{llama3.2:3b} nor \texttt{gemma3:4b} has reliable knowledge of ELTE-specific regulations — they were trained on general web text and have no exposure to the faculty's prerequisite rules or administrative procedures. The retrieval step injects the relevant passage directly into the prompt, allowing the model to produce a grounded answer even though it has no parametric knowledge of the topic.

\subsection{Limitations}

Several limitations should be acknowledged.

\paragraph{Small benchmark.} Twenty questions is a limited evaluation set. The in-scope questions are skewed toward prerequisite rules because that is the richest and most structured document in the corpus. Topics such as scholarship regulations, exam appeal procedures, and credit transfer rules are underrepresented.

\paragraph{Retrieval ceiling.} The 0.667 hit-rate means one third of in-scope questions are answered without relevant context. Improving chunking strategy (e.g.\ heading-aware splits, larger overlap) or adding a re-ranking step could raise this ceiling and correspondingly improve ROUGE-L and faithfulness.

\paragraph{ROUGE-L as a proxy.} ROUGE-L measures lexical overlap with a manually written reference answer. A response that is semantically correct but uses different phrasing will score poorly. For a more robust evaluation a human-judged correctness rating or a semantic similarity metric (e.g.\ BERTScore) would complement ROUGE-L.

\paragraph{Latency.} All configurations exceed the 30-second target. Achieving this target on CPU-only hardware with the current models would require either a smaller model (1B parameters) or a faster inference backend. For a deployment scenario where the server runs on shared faculty hardware with a GPU, latency would drop significantly.

\subsection{Recommended Configuration}

Based on the evaluation results, \textbf{gemma3:4b with RAG (B2)} is the recommended production configuration: it achieves the highest answer quality (ROUGE-L 0.534, faithfulness 0.844) and, counterintuitively, lower latency than the no-RAG Gemma baseline (61 s vs.\ 121 s). If latency is the primary constraint, \textbf{llama3.2:3b with RAG (B1)} offers a reasonable quality-speed trade-off at 70 seconds average response time with only a modest quality reduction.
